\documentclass[11pt]{article}

\usepackage{sectsty}
\usepackage{siunitx}
\usepackage{graphicx}
\usepackage{amsmath}
\usepackage{amsthm}
\usepackage{float}
\usepackage{outlines}
\usepackage{mathtools}
\usepackage[thinc]{esdiff}
\usepackage{amssymb}
\usepackage{outlines}
\usepackage{caption}
\usepackage{subcaption}
\let\oldsi\si
\renewcommand{\si}[1]{\oldsi[per-mode=reciprocal-positive-first]{#1}}
\usepackage{enumitem}
\newcommand{\degsym}{^{\circ}}
\newcommand{\Mod}[1]{\ (\mathrm{mod}\ #1)}
\usepackage{hyperref}
\hypersetup{
    colorlinks,
    citecolor=black,
    filecolor=black,
    linkcolor=black,
    urlcolor=black
}
\newcommand{\lb}{\\[8pt]}
\newenvironment*{cell}[1][]{\begin{tabular}[c]{@{}c@{}}}{\end{tabular}}
\newcommand{\img}[3]{\begin{center}
  \begin{figure}[H]
    \centering
    \includegraphics[width=#2\textwidth]{#1}
    \caption{#3}
    \label{fig:fig1}
  \end{figure}
\end{center}}
\newcommand{\doubleimg}[4]{\begin{center}
  \begin{figure}[H]
    \centering
    \begin{subfigure}{.45\textwidth}
      \centering
      \includegraphics[width=1\linewidth]{#1}
      \caption{#2}
      \label{fig:sub1}
    \end{subfigure}
    \begin{subfigure}{.45\textwidth}
      \centering
      \includegraphics[width=1\linewidth]{#3}
      \caption{#4}
      \label{fig:sub2}
    \end{subfigure}
  \end{figure}
\end{center}}
\newcommand{\dx}{\mathrm{~d}x}


\title{Math AA HL at KCA - Chapter 17 Notes}
\author{Tim Bao 2023-2025 Cohort}
\date{September 11th, 2024}

\begin{document}

\maketitle
\pagebreak
\tableofcontents
\pagebreak

\section{Basic Rules of Integration}

\begin{large}
  \begin{align*}
    \int u' \cdot u^n \dx & = \frac{u^{n+1}}{n+1} + C \\
    \int \frac{u'}{u} \dx & = \ln |u| + C             \\
    \int u' \cdot e^u \dx & = e^u + C                 \\
    \int u' \cdot a^u \dx & = \frac{a^u}{\ln a} + C
  \end{align*}
\end{large}

\pagebreak

\section{Reverse Chain Rule}
When you notice that a function is multiplied with its derivative, it's likely the reverse chain rule.

\pagebreak

\section{Trigonometric Integrals}

\begin{large}
  \addtolength{\jot}{1em}
  \begin{align*}
    \int \sin x \dx                      & = -\cos x + C                                                         \\
    \int \cos x \dx                      & = \sin x + C                                                          \\
    \int \tan x \dx                      & = -\ln |\cos x| + C                                                   \\ \\
    \int \sec^2 x \dx                    & = \tan x + C                                                          \\
    \int \sec x \tan x \dx               & = \sec x + C                                                          \\
    \int \csc^2 x \dx                    & = -\cot x + C                                                         \\
    \int \csc x \cot \dx                 & = -\csc x   + C                                                       \\ \\
    \int \frac{1}{\sqrt{a^2 - x^2}} \dx  & = \arcsin \left(\frac{x}{a}\right) + C                                \\
    \int -\frac{1}{\sqrt{a^2 - x^2}} \dx & = \arccos \left(\frac{x}{a}\right) + C                                \\
    \int \frac{1}{a^2 + x^2} \dx         & = \frac{1}{a} \arctan \left(\frac{x}{a}\right) + C \text{ with }x < a
  \end{align*}
\end{large}

\pagebreak
\section{Integration by Parts}
$$\int uv' \dx = uv - \int vu' \dx$$
\pagebreak

\section{Volume of Revolution}

If we revolve the curve $y = f(x)$ about the $x$-axis, the volume of the solid generated is given by $$V = \int_{a}^{b} \pi y^2 \dx$$




\end{document}